\documentclass[11pt]{amsart}
\usepackage[margin=1.25in]{geometry}
\usepackage{amsmath,amssymb,amsthm}
\usepackage{booktabs}
\usepackage[colorlinks=true,linkcolor=blue,citecolor=blue,urlcolor=blue]{hyperref}

\newtheorem{theorem}{Theorem}[section]
\newtheorem{proposition}[theorem]{Proposition}
\newtheorem{lemma}[theorem]{Lemma}
\theoremstyle{definition}
\newtheorem{definition}[theorem]{Definition}
\newtheorem{construction}[theorem]{Construction}
\theoremstyle{remark}
\newtheorem{remark}[theorem]{Remark}

\newcommand{\id}{\mathrm{id}}
\newcommand{\PSh}{\mathrm{PSh}}
\newcommand{\Sh}{\mathrm{Sh}}
\newcommand{\Set}{\mathbf{Set}}
\newcommand{\Noe}{\mathrm{Noetica}}
\newcommand{\Jleg}{J_{\mathrm{legacy}}}
\newcommand{\Eamb}{\mathcal{E}_{\mathrm{amb}}}
\newcommand{\Eres}{\mathcal{E}_{\mathrm{res}}}
\newcommand{\Eleg}{\mathcal{E}_{\mathrm{legacy}}}

\begin{document}

\title[The marked and modal Noetic topos]{Noetic Topos 2.0: a marked and modal successor to the conservative Noetic site}
\author{Little Henderson}
\thanks{White paper v0.1, 2026-08-07. Artifact status: \textsc{draft / gate 1 research candidate / ring 4 / not licensed / held}. Source basis: \texttt{NOETIC\_TOPOS\_2\_SPEC\_v0.1.md} (SHA-256 \texttt{156A5D8C\ldots}), theorem and proof-obligation ledgers v0.1, finite model certificate v0.1, modality-composition proof, publication delta Paper~2 at base commit \texttt{a1cd9901}. Drafting this paper promotes no claim.}
\date{August 7, 2026}

\begin{abstract}
The conservative Noetic Topos construction contains a valid conditional topos-existence argument, but its ACBE and RPM cover clauses do not retain the separate information later interpretations assigned to them: ACBE reflexivity places an identity in the corresponding sieve, and the RPM counit---split by its section---forces an identity by sieve closure. Each clause alone generates only maximal sieves, and the combined topology collapses to the jointly-surjective coverage already present in the source. Distinct modalities cannot be recovered by relabeling those generators.

Noetic Topos 2.0 is the first conservative successor that survives this collapse counterexample. It augments the legacy site with a finite four-atom context lattice $L=\mathcal{P}(\{0,1\}\times\{0,1\})$ whose coordinates carry formal ACBE and RPM marks, equips $L$ with two transverse coordinate-partition coverages, and keeps immutable, hash-chained cover provenance \emph{outside} the ordinary topology. On $\PSh(L,\Set)$, internal proofs backed by exhaustive finite checks establish that both generated topologies are subcanonical, their sheafification reflectors are distinct and non-identity, their join is the topology generated by the four-singleton partition, and the reflectors commute up to canonical natural isomorphism. The ambient candidate $\mathrm{NoeTop}_2=\PSh(\mathcal{C}\times L)$ carries three modal reflectors and a resolved subtopos; the full research object packages these with an external mark record. The product-site lift, legacy round-trip comparison, cross-context coupling, proof-assistant formalization of the sheaf layer, and independent category-theory review remain open; preserving those boundaries is treated as part of the result.
\end{abstract}

\maketitle

\section{The defect being repaired}

Let $\mathcal{C}$ denote the registered legacy category $\Noe$ and $\Jleg$ its proven coverage. The conservative source admits covering sieves from three routes: jointly-surjective observational families, designated ACBE descent maps, and the RPM counit. Two sieve arguments are decisive.

\begin{proposition}[Legacy collapse]\label{prop:collapse}
Under the registered axioms, the ACBE and RPM clauses contribute only maximal sieves, and hence
\[
\Jleg \;=\; J_{\text{jointly-surjective}} .
\]
\end{proposition}

\begin{proof}
(\emph{ACBE}) The reflexive descent map at equal contexts is $\id_X$. A sieve on $X$ containing $\id_X$ is maximal, because precomposition with every arrow into $X$ places every such arrow in the sieve.

(\emph{RPM}) The section-retraction data
\[
\eta_X : X \to R(X), \qquad \varepsilon_X : R(X) \to X, \qquad \varepsilon_X \circ \eta_X = \id_X
\]
imply that any sieve containing $\varepsilon_X$ also contains $\varepsilon_X \circ \eta_X = \id_X$ after closure under precomposition---again maximal.
\end{proof}

This does not invalidate the generic topos-existence statement; it narrows what the topology records: ordinary sieve generation has \emph{forgotten} the distinct ACBE and RPM provenance. A successor must change the cover generators; renaming the old maps is not a repair.

\begin{remark}[Subcanonicity gap]\label{rem:subcan}
There is no derivation of subcanonicity for the legacy topology. A finite countermodel (sets of size at most two with injections) exhibits a jointly-covering pair of point inclusions whose compatible local injections into a singleton admit no global injection. Until a revised site is proved subcanonical, the distinguished legacy object is represented in the sheaf topos by $a_{\mathrm{legacy}}(yX)$, never by silently treating $yX$ as a sheaf \cite{stacks00WO}.
\end{remark}

\section{The finite marked context factor}

\begin{definition}[Context lattice]\label{def:lattice}
Let
\[
\Omega = \{0,1\}\times\{0,1\}, \qquad L = \mathcal{P}(\Omega), \qquad (U \to V \iff U \subseteq V).
\]
The first coordinate is the formal ACBE mark and the second the formal RPM mark. These are structural labels only, with no empirical interpretation until a separate operationalization and validation record exists.
\end{definition}

\begin{definition}[Transverse coverages]\label{def:covers}
With blocks $A_i=\{(i,0),(i,1)\}$ and $R_j=\{(0,j),(1,j)\}$, the generating families at each context $V \in L$ are
\begin{align*}
K_A(V) &= \{\, V\cap A_0 \to V,\;\; V\cap A_1 \to V \,\}, \\
K_R(V) &= \{\, V\cap R_0 \to V,\;\; V\cap R_1 \to V \,\},
\end{align*}
with the empty family covering the empty context.
\end{definition}

Unlike the legacy clauses, the global pieces are \emph{proper} subcontexts: neither family contains an identity. Pullback along $U \to V$ is intersection with $U$, so each basis is stable under pullback; refining a partition piece by the same partition returns that piece plus empty pieces.

\begin{construction}[Marks as parallel audit data]\label{con:marks}
An ordinary Grothendieck topology retains generated sieves, not why a family was proposed. The reference implementation therefore maintains two representations. The topology consumes the unmarked family; governance code consumes an immutable marked cover
\[
\bigl(\text{derivation kind},\ \text{target},\ \text{exact member arrows},\ \text{mark record},\ \text{parent digests}\bigr),
\]
where the derivation kind is \textsc{generator}, \textsc{pullback}, or \textsc{common refinement}, and the mark record binds cover kind, source authority, assumptions, evidence references, review state, and provenance ancestry. Canonical JSON and SHA-256 bind the target, members, mark, and sorted parent digests. Deserialization of a derived record resolves and rehashes every parent, checks the acyclic derivation, and recomputes child members and mark; a common refinement without exactly one ACBE and one RPM generating parent fails closed. At singleton contexts the bare ACBE and RPM families coincide, but their marked records and digests remain distinct.
\end{construction}

\begin{remark}
The mark mechanism is an external serialization and audit invariant---not a claim that a marked cover is an internal topos object, a colored-site construction, or a categorical semantics of evidence. Conflating the two interfaces would repeat the original information-loss error.
\end{remark}

\section{The ambient modal candidate}

The successor is \emph{not} defined as the already-joined sheaf category: that would make the ACBE and RPM reflectors identity on every admitted object and erase the modal structure a second time. Instead, begin with the ambient presheaf topos and three topology refinements:
\[
\Eamb = \PSh(\mathcal{C}\times L), \qquad
J_{\mathrm{res}} = J_{\mathrm{obs}} \vee J_A \vee J_R, \qquad
\Eres = \Sh(\mathcal{C}\times L,\ J_{\mathrm{res}}),
\]
with sheafification-then-inclusion reflectors $M_{\mathrm{obs}}=i_{\mathrm{obs}}a_{\mathrm{obs}}$, $M_A=i_A a_A$, $M_R=i_R a_R$ \cite{stacks00W1}. Assuming $\mathcal{C}$ is small in the chosen universe,
\[
\mathrm{NoeTop}_2 := \Eamb, \qquad
\mathrm{NoeSystem}_2 = \bigl(\mathrm{NoeTop}_2,\ M_{\mathrm{obs}},\ M_A,\ M_R,\ \Eres,\ \text{external mark record}\bigr).
\]
$\mathrm{NoeSystem}_2$ is a structured package, not itself a topos or category; the provenance record is external audit data. Raw and partially resolved evidence lives in the ambient topos; only jointly resolved objects enter $\Eres$. On $\Eres$, both $M_A$ and $M_R$ restrict to identity---no text may describe them as nontrivial endomodalities of the resolved subtopos.

The v0.1 executable model covers only the $L$ coordinate and its mark derivations. Typed observational covers in $\mathcal{C}$, general arrows of $\mathcal{C}\times L$, and the product-site comparison remain open.

\section{Theorems on the finite context factor}

\begin{theorem}[Generated topologies]\label{thm:nt2-1}
Each of $K_A$, $K_R$ is jointly exhaustive, stable under pullback, and stable under its own cover composition; the generated covering sieves form Grothendieck topologies $J_A$, $J_R$ on $L$.
\end{theorem}

\begin{proof}
Pullbacks in $L$ are intersections. Pulling $V\cap A_i \to V$ back along $U \to V$ gives $U\cap A_i \to U$, and similarly for RPM. The union of the two pieces is $V$. Refining either piece by the same coordinate partition returns that piece together with an empty piece, so the composite family refines the original cover. Identity families and the null cover at the empty context give the remaining pretopology cases; the standard passage from a pretopology to covering sieves \cite{stacks00ZC} concludes.
\end{proof}

\begin{theorem}[Subcanonicity of the context factor]\label{thm:nt2-2}
Both context-factor topologies are subcanonical on $L$.
\end{theorem}

\begin{proof}
Let $yW$ be representable. A section of $yW$ on $V$ exists exactly when $V \subseteq W$, and is then unique. If both pieces of a marked cover map into $W$, their union $V$ maps into $W$; conversely a map from $V$ restricts uniquely. The empty case is also unique. Thus the representable gluing map is a bijection for every generator.
\end{proof}

This theorem applies only to $L$; it does not repair the legacy subcanonicity gap of Remark~\ref{rem:subcan}.

\begin{theorem}[Distinct, non-identity modalities]\label{thm:nt2-3}
The ACBE and RPM sheafification modalities on $\PSh(L,\Set)$ are distinct and neither is the identity modality.
\end{theorem}

\begin{proof}
Define $F_A(V)$ to be the set of bit assignments to those complete ACBE blocks $A_i$ contained in $V$, with restriction forgetting missing blocks. Complete ACBE blocks split uniquely across $K_A(V)$, so $F_A$ satisfies ACBE gluing. At $V=\Omega$, $F_A(\Omega)$ has four sections, while each RPM half contains no complete ACBE block; the RPM matching-family set has one element, so the RPM gluing map is not injective. Thus $F_A$ is a $J_A$-sheaf but not a $J_R$-sheaf. The symmetric construction $F_R$ reverses the roles. Therefore the reflective subcategories and their reflectors differ, and each reflector changes at least one presheaf.
\end{proof}

\begin{remark}[Ambient lift caveat]
An earlier external-product distinctness witness for $\PSh(\mathcal{C}\times L)$ is false in general: at the null cover over $(c,\varnothing)$, $y\mathcal{C}_0(c)\times F_A(\varnothing)$ need not be a singleton, and duplicated representable coordinates can obstruct gluing. Ambient distinctness therefore remains open; the valid result is limited to $\PSh(L,\Set)$.
\end{remark}

\begin{theorem}[Joined topology]\label{thm:nt2-4}
The common refinement of the ACBE and RPM global covers is the four-singleton cover, and
\[
J_A \vee J_R \;=\; J_{\mathrm{singleton}},
\]
the topology generated by the singleton partition together with the null cover at the empty context.
\end{theorem}

\begin{proof}
Every intersection $A_i\cap R_j$ is the singleton $\{(i,j)\}$. At an arbitrary $V$, the refinement consists of the singletons belonging to $V$: the atomic cover. Every jointly exhaustive family of subcontexts contains, in its generated sieve, each singleton of $V$; conversely every generator is jointly exhaustive. Thus the joined context coverage is the ordinary union topology on the finite powerset lattice.
\end{proof}

The internal label $J_{\mathrm{atom}}$ used by the source artifacts does not assert the standard category-theoretic atomic topology; $J_{\mathrm{singleton}}$ is the preferred name at the publication layer.

\begin{theorem}[Modality composition]\label{thm:composition}
For the explicit reflectors
\[
(R_P F)(V) \;=\; \prod_{\substack{B \in P \\ V\cap B \neq \varnothing}} F(V\cap B)
\qquad (P \in \{A, R, \mathrm{atom}\}),
\]
with the empty product a singleton, each $R_P$ realizes the associated-sheaf reflector for $J_P$ in this finite setting, and there are canonical natural isomorphisms on $\PSh(L,\Set)$:
\[
M_A M_R \;\cong\; M_{\mathrm{atom}} \;\cong\; M_R M_A .
\]
\end{theorem}

\begin{proof}[Proof sketch]
The registered proof checks the universal property directly: $R_P F$ is a $J_P$-sheaf, and every map from $F$ to a $J_P$-sheaf factors uniquely through the unit. Expanding the two transverse products,
\[
(R_A R_R F)(V) \;=\; \prod_i \prod_j F(V\cap A_i\cap R_j) \;\cong\; \prod_{\omega\in V} F(\{\omega\}) \;=\; (R_{\mathrm{atom}}F)(V),
\]
naturally in $F$ and $V$; the inverse groups atomic coordinates by blocks, and swapping $A$ and $R$ gives the other direction. Both comparisons respect restrictions, natural transformations, and reflector units. These are natural isomorphisms, not strict equalities.
\end{proof}

\begin{remark}[Distinct but commuting]
Theorem~\ref{thm:composition} corrects an important narrative risk: the finite modalities are distinct and non-identity (Theorem~\ref{thm:nt2-3}) but they \emph{commute}. Distinctness does not imply noncommutation, and no current artifact licenses a noncommutation claim.
\end{remark}

\subsection*{Executable corroboration}
A dependency-free checker exhausts all sixteen contexts, all sixteen representables, both witness presheaves, and every generating-cover pullback. The exact-topology implementation enumerates all 298 sieves over the sixteen contexts and reports 168 covering sieves for both the join and singleton presentations; the modality verifier covers 18 fixtures and 3{,}012 restriction squares. At the bound commit the finite-site suite (8 tests), the modality-composition suite (6 tests), and the pinned Lean Gate~1 verifier all pass. The Lean scope covers four-atom exhaustiveness, proper global pieces, transverse refinement, and witness separation---not the sheaf conditions, topology generation, reflector universal properties, or product-site comparison. Execution is corroboration of finite proofs, not formal verification or independent review.

\section{Gluing boundary for probabilistic data}

For $\Set$-valued presheaves, gluing is bijectivity of the canonical map from a global section to matching local sections. Probability distributions do not inherit this by analogy: pairwise-compatible marginals may lack a global joint distribution. A probabilistic successor must pass one of four explicit gates: an acyclic or junction-tree cover with a proved extension theorem; a finite linear-program feasibility certificate for the joint distribution; an obstruction calculation with a stated coefficient system; or a quantitative relaxation whose residual and uncertainty are reported. No passing $\Set$-valued sheaf test licenses a probabilistic consistency claim.

\section{Backward compatibility as a conjectural comparison}

The legacy topos is preserved verbatim as $\Eleg=\Sh(\mathcal{C},\Jleg)$. The target comparison is
\[
\Eres \;\cong\; \Eleg^{\,4}
\]
after atomic resolution of the context factor, with a proposed fully faithful neutral-channel embedding $X \mapsto (X,1,1,1)$ and a specified natural isomorphism $\mathrm{ev}_{00}\circ\iota \cong \mathrm{Id}_{\Eleg}$. A categorical natural isomorphism is not a bit-for-bit serialization receipt: identifier and canonical-payload-hash preservation is a separate strict-equality requirement. Both round trips remain proof-required; this paper claims no backward compatibility.

Even if proved, $\Eleg^{4}$ exposes a strict expressivity boundary: four independent legacy channels with no remaining overlap, transition law, or cross-atom coupling. The sixteen-context construction is a finite separation and provenance reference model---not yet a substantively rich final semantics. A final successor must either name the ambient modal package as its principal semantics or add typed overlaps and transition data and prove their coherence; which route is canonical remains open, and v0.1 licenses neither.

\section{Admission gates}

Noetic Topos 2.0 may leave Ring~4 only after all of the following are present: (1) independent category-theory review of Proposition~\ref{prop:collapse} and Theorems~\ref{thm:nt2-1}--\ref{thm:nt2-4}; (2) proof-assistant artifacts for the finite site, witness presheaves, and product-site comparison; (3) a verified legacy round trip with stable identifiers and hashes; (4) a theorem-ledger merge linking every statement to exact source and proof artifacts; (5) proof and review of the product-site lift of Theorem~\ref{thm:composition} and its interaction with observational coverage; (6) separate probability-gluing theorems or feasibility certificates; (7) a canonical decision and proof package for ambient semantics versus genuine cross-context coupling; and (8) explicit confirmation that no empirical or product wording was licensed.

\section{Conclusion}

The correction and the successor fit together cleanly: the legacy ACBE/RPM clauses collapse because identities enter their sieves, and nontriviality is recovered only by changing the context generators. In the finite model the new modalities are distinct and non-identity, but they commute and jointly resolve to four atomic channels, while hash-chained mark records preserve the derivational provenance ordinary topology generation forgets. That is a coherent mathematical reference point, not a finished semantics of contextual cognition---and maintaining that distinction is itself part of the contribution.

\begin{thebibliography}{9}
\bibitem{stacks00ZC} The Stacks Project Authors, ``The topology defined by a site,'' Tag 00ZC, \url{https://stacks.math.columbia.edu/tag/00ZC}.
\bibitem{stacks00W1} The Stacks Project Authors, ``Sheafification,'' Tag 00W1, \url{https://stacks.math.columbia.edu/tag/00W1}.
\bibitem{stacks00WO} The Stacks Project Authors, ``Representable sheaves,'' Tag 00WO, \url{https://stacks.math.columbia.edu/tag/00WO}.
\bibitem{maclane} S.~Mac Lane and I.~Moerdijk, \emph{Sheaves in Geometry and Logic}, Springer, 1994.
\bibitem{redteam} Project Coherence, \emph{Mathematical Red-Team Report}, v0.1 (collapse counterexample and subcanonicity countermodel).
\end{thebibliography}

\end{document}
